\section{Value Function：情感作为人类的奖励信号}

\subsection{当前 RL 的局限}

当前的 RL 训练存在一个根本问题：模型需要完成整个轨迹才能获得奖励信号。如果任务很长（数百到数千步），在得到最终结果之前不会有任何学习发生。

\begin{importantbox}{Value Function 的作用}
Value function 允许在中间状态就判断"你做得好不好"。类比下棋：丢掉一个棋子时你就知道自己搞砸了，不需要把整盘棋下完。同样，如果你在思考数学问题时花了 1000 步探索一个方向，最后发现不对，value function 可以在那一刻就发出信号——"1000 步之前你就不该走这条路"。
\end{importantbox}

\subsection{情感作为 Value Function}

\begin{importantbox}{一个令人深思的案例}
一个因脑损伤失去情感处理能力的人：他仍然口齿清晰，能解小谜题，测试成绩正常——但他变得极其不擅长做任何决策。穿哪双袜子要花几个小时，还做出很糟糕的财务决策。

这说明：\textbf{情感是人类决策的核心价值函数}。没有情感，即使"智力"完好，行为能力也严重退化。
\end{importantbox}

\begin{knowledgebox}{情感的简单性与鲁棒性的权衡}
情感可能足够简单以至于可以被完整描述（"map them out in a human understandable way"）。但正因为简单，它们才在完全不同于进化环境的现代世界中仍然有效——这是复杂性与鲁棒性之间的经典权衡。不过也有反例：人类的饥饿感在食物过剩的现代世界中就不那么准确了。
\end{knowledgebox}

\subsection{本章小结}

Value function 将使 RL 训练更高效，但 Ilya 认为它不是根本性的突破——"anything you can do with a value function you can do without, just more slowly"。真正关键的是泛化问题。情感作为人类的 value function，其可靠性和简单性值得 AI 研究者深思。


